\documentclass[11pt,a4paper]{article}
\usepackage[utf8]{inputenc}
\usepackage[T1]{fontenc}
\usepackage{geometry}
\geometry{a4paper, margin=1in}
\usepackage{xcolor}
\usepackage{hyperref}
\usepackage{titlesec}
\usepackage{enumitem}
\usepackage{tikz}
\usetikzlibrary{shapes.geometric, arrows.meta, positioning, calc}
\usepackage{tcolorbox}
\usepackage{booktabs}
\usepackage[french]{babel}

% Couleurs
\definecolor{riotred}{HTML}{FF4655}
\definecolor{tacticalcyan}{HTML}{00F5D4}
\definecolor{neuralblue}{HTML}{3B82F6}
\definecolor{darkbg}{HTML}{020617}

% Style des sections
\titleformat{\section}{\large\bfseries\sffamily\color{neuralblue}}{\thesection}{1em}{}
\titleformat{\subsection}{\normalsize\bfseries\sffamily\color{neuralblue}}{\thesubsection}{1em}{}

% Configuration hyperref
\hypersetup{
    colorlinks=true,
    linkcolor=neuralblue,
    urlcolor=neuralblue,
}

\title{\sffamily\bfseries\huge \color{neuralblue} AURORA // MAP AGENT WORKFLOW}
\author{Système d'Intelligence AURORA}
\date{\today}

\begin{document}

\maketitle

\vspace{-1em}
\begin{center}
    \textit{Module d'Intelligence Spatiale pour la Vigilance Tactique}
\end{center}
\vspace{1em}

\noindent L'\textbf{Map Agent} est un module d'intelligence spatiale spécialisé au sein de l'écosystème tactique AURORA :
\begin{itemize}
    \item Il exploite les LLM multi-modaux (\textbf{Gemini Vision}).
    \item Il utilise les métadonnées officielles de \textbf{Riot Games}.
    \item Il fournit une conscience géographique et tactique en temps réel à partir des visuels de jeu.
\end{itemize}

\section{Description Générale \& Objectif}
L'Map Agent transforme les \textbf{visuels bruts} en \textbf{intelligence tactique structurée} via une approche hybride :
\begin{itemize}[label=\textbullet]
    \item \textbf{Visual Recognition} : Identification de la carte via Gemini Vision.
    \item \textbf{Data Grounding} : Validation par l'API officielle de Riot Games.
    \item \textbf{Lore Synthesis} : Enrichissement narratif et géographique.
\end{itemize}

\textbf{Objectifs Clés :}
\begin{itemize}[label=$\diamond$]
    \item Automatisation de la détection spatiale (carte et callouts).
    \item Calcul de coordonnées mondiales précises.
    \item Diffusion de briefings tactiques et narratifs.
\end{itemize}

\section{System Architecture}
Le flux de travail suit un pipeline linéaire allant de l'\textbf{Ingestion} à l'\textbf{Enrichissement}.

\begin{center}
\begin{tikzpicture}[
    node distance=1.6cm and 2.2cm,
    block/.style={rectangle, draw, rounded corners, minimum width=3cm, minimum height=0.8cm, text centered, font=\sffamily\scriptsize},
    io/.style={ellipse, draw, fill=gray!5, minimum width=2.8cm, minimum height=0.8cm, text centered, font=\sffamily\scriptsize},
    arrow/.style={-{Stealth[scale=1.0]}, rounded corners}
]
    % Central Column
    \node [io] (user) {User Screenshot};
    \node [block, below=of user] (frontend) {Frontend Interface};
    \node [block, below=of frontend] (server) {Map Agent Server};
    \node [block, below=of server, yshift=-0.5cm] (gemini) {Gemini Vision};
    \node [block, below=of gemini] (oracle) {Metadata Oracle};
    \node [block, below=of oracle] (lore) {Lore Geographer};
    \node [io, below=of lore] (report) {Final Intel Report};

    % Side Nodes
    \node [block, left=of gemini] (ocr) {Minimap OCR};
    \node [block, right=of gemini] (tracer) {Precision Tracer};
    \node [block, right=of oracle] (db) {Riot Data};

    % Main Vertical Flow
    \draw [arrow] (user) -- (frontend);
    \draw [arrow] (frontend) -- (server);
    \draw [arrow] (server) -- (gemini);
    \draw [arrow] (gemini) -- (oracle);
    \draw [arrow] (oracle) -- (lore);
    \draw [arrow] (lore) -- (report);

    % Sideways Distribution from Server
    \draw [arrow] (server.west) -| (ocr.north);
    \draw [arrow] (server.east) -| (tracer.north);
    
    % Data & Processing
    \draw [arrow] (oracle) -- (db);
    
    % Convergence to Report
    \draw [arrow] (ocr.south) |- (report.west);
    \draw [arrow] (tracer.south) |- (report.east);
    
    % Feedback Loop
    \draw [arrow] (report.west) -- ++(-1.5,0) |- (frontend.west);

    % Core Intelligence Box
    \draw[dashed, neuralblue!40] ($(ocr.north west)+(-0.4,0.4)$ ) rectangle ($(db.south east)+(0.4,-0.4)$);
    \node[anchor=north east, neuralblue, font=\sffamily\tiny] at ($(db.south east)+(0.4,-0.4)$) {Core Intelligence};
    
\end{tikzpicture}
\end{center}

\section{Analyse des Composants}

\subsection{Map Agent Server (\texttt{map\_agent\_server.py})}
L'orchestrateur central de l'ensemble du flux de travail. Il gère le cycle de vie d'une demande d'analyse.
\begin{itemize}
    \item \textbf{Input} : Flux d'octets d'image brute (PNG/JPG).
    \item \textbf{Process} : Effectue l'OCR pré-analyse ; Synchronise les appels ; Exécute le traçage de précision.
    \item \textbf{Output} : Objet d'intelligence tactique JSON unifié.
\end{itemize}

\subsection{Metadata Oracle (\texttt{metadata\_oracle.py})}
La source faisant autorité qui convertit le langage naturel en données spatiales.
\begin{itemize}
    \item \textbf{Input} : Nom de carte prédit et chaîne d'emplacement.
    \item \textbf{Process} : Correspondance floue avec les données de Riot ; Calcule les niveaux d'élévation.
    \item \textbf{Output} : Coordonnées officielles (X, Y, Z), Super-région et Conseils Tactiques.
\end{itemize}

\subsection{Lore Geographer (\texttt{lore\_geographer.py})}
Ajoute une couche de narration et de contexte réel aux données tactiques.
\begin{itemize}
    \item \textbf{Input} : Lore Coordinates (Lat/Long).
    \item \textbf{Process} : Associe les coordonnées à des lieux réels ; Génère des extraits de briefing.
    \item \textbf{Output} : Nom du lieu réel et « Tactical Lore Report ».
\end{itemize}

\section{Résumé du Flux de Données}

\begin{center}
\small
\begin{tabular}{@{}llll@{}}
\toprule
\textbf{Composant} & \textbf{Input} & \textbf{Primary Tool} & \textbf{Output} \\ \midrule
\textbf{Server}    & Screenshot     & FastAPI / Gemini      & Raw Map ID      \\
\textbf{Oracle}     & Raw Map ID     & Fuzzy Matching        & Exact Coords    \\
\textbf{Geographer}  & Lore Coords    & Gemini NLP            & Geo-Context     \\
\textbf{Tracer}    & Minimap Pixel  & Matrix Math           & World Coords \\ \bottomrule
\end{tabular}
\end{center}

\begin{tcolorbox}[colback=tacticalcyan!5,colframe=tacticalcyan!80,title=\sffamily\bfseries TIP : Precision Tracing]
\small
\textbf{Precision Tracing} : Localisation 3D via icône minimap et multiplicateurs Riot officiels.
\end{tcolorbox}

\end{document}
